%-------------------------
% Cover letter -- Genmab, Translational Data Science internship: hierarchical
% and Bayesian modelling of historical control-arm data from preclinical
% tumour studies (Early Translational Data Science team, Utrecht). Lives on
% this branch only.
%
% Same layout as the ASML letter: the resume's header, 1in sides, 0.75in top
% and bottom. Build: cd letters; pdflatex genmab-cover-letter.tex
%
% Project claims come from cv/resume.tex on this branch; the Biohub
% validation detail comes from the Volt-biohub commit log; the coursework
% sentence is as Alex stated it (Oct 2026). No R anywhere: the posting asks
% for "R and/or Python", so Python covers it.
% Style rules for this letter: plain words, no em-dashes, no "not X but Y".
%-------------------------

\documentclass[letterpaper,11pt]{article}

\usepackage[left=1in,right=1in,top=0.75in,bottom=0.75in]{geometry}
\usepackage[hidelinks]{hyperref}
\usepackage[english]{babel}

\urlstyle{same}
\pagestyle{empty}
\raggedright
\setlength{\parindent}{0pt}
\setlength{\parskip}{6pt}

% Same guard as resume.tex: machine-readable text under pdflatex, still
% compiles under XeTeX/tectonic.
\ifdefined\pdfgentounicode
  \input{glyphtounicode}
  \pdfgentounicode=1
\fi

\begin{document}

%-----------HEADING (identical to resume.tex)-----------
\begin{center}
    \textbf{\Huge \scshape Oleksii Krasnoshtanov} \\ \vspace{1pt}
    \small
    \href{mailto:oleksiikrasnoshtanov@gmail.com}{\underline{oleksiikrasnoshtanov@gmail.com}} $|$
    +31 6 29746862 $|$
    Tilburg, Netherlands \\ \vspace{2pt}
    \href{https://www.linkedin.com/in/oleksii-krasnoshtanov/}{\underline{linkedin.com/in/oleksii-krasnoshtanov}} $|$
    \href{https://github.com/alex-krasnoshtanov}{\underline{github.com/alex-krasnoshtanov}}
\end{center}
% The resume's section rule, reused as the header divider.
\vspace{-4pt}\noindent\rule{\textwidth}{0.4pt}

\today

Genmab, Early Translational Data Science \\
Utrecht, Netherlands

% Replace with the recruiter's name if the portal shows one.
Dear Hiring Manager,

% Genmab posts several projects in this department and allows two
% applications. If you apply to a second one, both letters must name the
% project titles in order of preference, with a line on each.
I am applying for the internship in the Early Translational Data Science team
on modelling historical control-arm data from preclinical tumour studies. I am
a third-year student in Applied Data Science and Artificial Intelligence at
Breda University of Applied Sciences, where I study on a scholarship with a
9.1/10 average.

Genmab appeals to me because its work stays close to both research and
patients. Contributing to medicine in any form matters a lot to me, and working
on data behind cancer treatment would be very motivating. I also like that this
project has a clear practical aim: if historical controls can be reused safely,
future studies can use fewer animals and still keep their statistical power.

% Biohub facts (Volt-biohub commit log): validation by leaving one embryo out;
% same-embryo hold-outs mis-ranked variants, a cross-embryo check agreed with
% the board; noise floor measured; division knobs flipped sign by embryo.
% "Top 14%": #542 of ~3,970 on the private leaderboard.
The closest work I have done is a Kaggle competition on tracking cells in 3D
microscopy of developing embryos, where my team of two finished in the top 14\%
of about 4,000 teams. The data came from only a few embryos, so we validated by
leaving one embryo out. When hold-outs from embryos the model had already seen
started ranking our changes wrongly, we built a cross-embryo check that agreed
with the leaderboard. We also measured the noise floor of our scores and
dropped settings whose effect changed sign from one embryo to another. This
project asks a similar question in a more structured setting: how much of the
variation comes from the animal and how much from the study, and whether a
model still holds when one study is left out.

Statistical modelling, including Bayesian methods, was part of my first year,
and in later projects we used statistical models as the baseline before moving
on to machine learning. I would like to take this further into hierarchical
models, which is a large part of why this project interests me. I have also
worked with biological data before. I ranked 2nd of 77 in a Kaggle competition
on plant root segmentation for NPEC, and in a five-person team I worked on a
plant phenotyping platform where researcher feedback on flagged predictions
drove retraining.

I enjoy working with people from other fields. In my current project with CEVA
Logistics and bol, I work with several companies to get access to their data,
help write business requirements documents, and make sure our proposals are
clear to people without a technical background. I ask a lot of questions when I
am new to a field, and I expect to do that often with the scientists in
Non-Clinical Safety and Pre-clinical in vivo Pharmacology.

% Availability: the posting requires 9 months (minimum 8), full-time, starting
% 1 February or 1 March. The last availability Alex stated was six months from
% early 2027, so no duration is claimed here. Add a sentence once confirmed.
I work mainly in Python, with Pandas and NumPy for data work. I use Git for
everything and keep my code tested and documented; the platform project ran an
85\% coverage gate across 430 tests. Utrecht is a direct train ride from
Tilburg, so three days a week on site works well for me.

I would be glad to talk about how I could contribute to the project. Thank you
for your time and consideration.

Kind regards,

Oleksii Krasnoshtanov

\end{document}
